\chapter{Cahier des Charges}

\begin{spacing}{1.25}
{\large

\section{Présentation de l'organisme d'accueil}

\subsection{N+ONE DATACENTERS}

\todo[inline, color=orange!30]{À RÉDIGER : Présentation N+ONE DATACENTERS — date de création, siège social (Casablanca), activité principale (hébergement/gestion datacenter), nombre de datacenters et capacité, clients types (entreprises, opérateurs, administrations), effectif approximatif, certifications (ISO 27001, Tier III ?)}

\subsection{Environnement de travail}

\todo[inline, color=orange!30]{À COMPLÉTER : décrire l'équipe, le département d'accueil, les outils utilisés en interne}

Le stage s'est déroulé au sein du département \textit{\colorbox{yellow}{[NOM DU DÉPARTEMENT]}},
sous la supervision de M. Ayoub \textsc{Amellaln} et M. Amine \textsc{El Mahi}.

\begin{figure}[ht]
\centering
\fbox{%
  \begin{minipage}{0.85\textwidth}
    \vspace{2cm}
    \centering
    \textit{[À INSÉRER — Organigramme de N+ONE DATACENTERS — images/organigramme.png]}
    \vspace{2cm}
  \end{minipage}%
}
\caption{Organigramme de N+ONE DATACENTERS}
\label{fig:organigramme}
\end{figure}

% -------------------------------------------------------------------

\section{Présentation du sujet}

\subsection{Les appels d'offres}

Un appel d'offres est une procédure formalisée par laquelle une organisation (le
\textit{maître d'ouvrage}) invite des fournisseurs ou prestataires potentiels à
soumettre des propositions commerciales et techniques pour la réalisation d'un marché.
Cette procédure garantit la mise en concurrence, la transparence et la traçabilité
des décisions d'achat.

Le cycle de vie d'un appel d'offres se décompose en plusieurs phases distinctes,
illustrées dans la figure~\ref{fig:lifecycle} :

\begin{enumerate}
    \item \textbf{BROUILLON (DRAFT)} : l'offre est en cours de rédaction,
    non visible par les candidats.
    \item \textbf{PUBLIÉE (PUBLISHED)} : l'offre est ouverte, les sociétés
    externes peuvent soumettre leurs propositions via le portail public.
    \item \textbf{ÉVALUATION} : la date limite de soumission est passée ;
    les responsables évaluent et scorent les candidatures reçues.
    \item \textbf{ATTRIBUÉE (AWARDED)} : le marché est attribué au candidat retenu.
    \item \textbf{ANNULÉE (CANCELLED)} : l'offre est interrompue à tout moment
    avant attribution.
\end{enumerate}

\begin{figure}[ht]
\centering
\fbox{%
  \begin{minipage}{0.85\textwidth}
    \vspace{2cm}
    \centering
    \textit{[À INSÉRER — Cycle de vie d'un appel d'offres — images/lifecycle.png]}
    \vspace{2cm}
  \end{minipage}%
}
\caption{Cycle de vie d'un appel d'offres}
\label{fig:lifecycle}
\end{figure}

\subsection{Les acteurs du système}

Le système distingue deux catégories d'utilisateurs :

\paragraph{Acteurs internes} Membres du personnel de N+ONE disposant d'un compte
applicatif, répartis en quatre rôles :

\begin{itemize}
    \item \textbf{ADMIN} : administrateur système, gère les comptes utilisateurs
    et accède au journal d'audit complet.
    \item \textbf{MANAGER} : responsable d'offres, crée et pilote les appels d'offres,
    attribue les marchés, génère les rapports PDF.
    \item \textbf{ANALYST} : évaluateur, score les soumissions et consulte
    les dossiers des candidats.
    \item \textbf{VIEWER} : lecteur seul, consulte les offres publiées
    sans pouvoir modifier quoi que ce soit.
\end{itemize}

\paragraph{Acteurs externes} Sociétés soumissionnaires accédant uniquement au
portail public sans créer de compte. Elles peuvent consulter les offres ouvertes
et déposer leurs propositions directement en ligne.

% -------------------------------------------------------------------

\section{Analyse des besoins}

\subsection{Besoins fonctionnels}

\begin{table}[ht]
\centering
\caption{Besoins fonctionnels du système}
\label{tab:besoins-fonctionnels}
\begin{tabularx}{\textwidth}{|l|X|}
\hline
\textbf{Module} & \textbf{Description} \\
\hline
Authentification & Connexion sécurisée, déconnexion, rafraîchissement de session \\
\hline
Gestion des offres & Création, modification, publication, transition de statut, annulation \\
\hline
Candidats et soumissions & Enregistrement des candidats, dépôt de soumissions, scoring technique et financier \\
\hline
Documents & Upload multipart, téléchargement sécurisé, chiffrement sélectif \\
\hline
Portail public & Consultation des offres publiées et dépôt de soumissions sans authentification \\
\hline
Notifications & Alertes automatiques : délais approchants, changements de statut, attribution \\
\hline
Rapports & Export PDF d'une offre (synthèse, tableau des scores, candidat retenu) \\
\hline
Audit & Journal des actions sur toutes les ressources (qui, quoi, quand, depuis où) \\
\hline
Administration & Gestion des comptes utilisateurs (création, désactivation, changement de rôle) \\
\hline
\end{tabularx}
\end{table}

\subsection{Besoins non-fonctionnels}

\begin{description}
    \item[Sécurité] Authentification par JWT (token en mémoire, cookie httpOnly),
    contrôle d'accès basé sur les rôles (RBAC), chiffrement AES-256-GCM des documents
    sensibles, audit complet, protection contre les injections SQL, HTTPS obligatoire.

    \item[Performance] Réponses API inférieures à 500 ms pour les listes paginées.
    Pagination côté serveur (20 offres par page). Chiffrement uniquement sur les
    documents marqués sensibles pour ne pas pénaliser les uploads standards.

    \item[Disponibilité] Redémarrage automatique des conteneurs Docker en cas d'échec
    (\texttt{restart: unless-stopped}). Vérifications de santé (\textit{health checks})
    sur chaque service.

    \item[Traçabilité] Toute action d'écriture (création, modification, téléchargement,
    export, connexion) est enregistrée dans le journal d'audit avec l'adresse IP,
    l'utilisateur et le contenu des modifications.

    \item[Maintenabilité] Architecture modulaire (8 applications Django indépendantes),
    pipeline CI automatique (tests + lint à chaque push), structure de projet
    documentée.

    \item[Reproductibilité] Déploiement en une commande (\texttt{docker-compose up}).
    Variables d'environnement pour tous les secrets, aucune valeur en dur dans le code.
\end{description}

% -------------------------------------------------------------------

\section{Évaluation et justification de la solution}

\subsection{Étude de l'existant}

Plusieurs solutions existent sur le marché pour la gestion des appels d'offres :

\begin{table}[ht]
\centering
\caption{Comparaison des solutions existantes}
\label{tab:existant}
\begin{tabularx}{\textwidth}{|l|X|X|X|}
\hline
\textbf{Solution} & \textbf{Points forts} & \textbf{Limites} & \textbf{Adapté ?} \\
\hline
SAP Ariba & Complète, intégration ERP & Coût élevé (licences), délai de déploiement long, sur-mesure impossible & Non \\
\hline
Proactis & Gestion marchés publics & Coût prohibitif pour une PME, rigidité du workflow & Non \\
\hline
OpenProposal / open source & Gratuit & Documentation insuffisante, peu maintenu, pas de support & Non \\
\hline
Excel + partage réseau & Gratuit, familier & Non sécurisé, non traçable, fragile, pas de portail externe & Non \\
\hline
\textbf{Développement sur mesure} & \textbf{Adapté exactement aux besoins} & Délai de développement & \textbf{Oui} \\
\hline
\end{tabularx}
\end{table}

\subsection{Critique de l'existant}

Aucune des solutions commerciales n'est adaptée au contexte de ce stage pour
trois raisons cumulatives. D'abord, le \textbf{budget nul} interdit les licences SaaS.
Ensuite, le \textbf{délai d'un mois} exclut tout déploiement et personnalisation
de produit tiers. Enfin, le besoin d'un \textbf{chiffrement sélectif} des documents
financiers et d'un portail public isolé dépasse les fonctionnalités standard
des outils génériques.

La solution Excel actuelle, bien que pratique, expose l'entreprise à des risques
réels : perte de données, accès non contrôlé aux offres financières, impossibilité
de retrouver l'historique des modifications.

\subsection{Justification du développement sur mesure}

Le développement d'une application web sur mesure est la seule option répondant
simultanément à toutes les contraintes :

\begin{itemize}
    \item Workflow exact du processus N+ONE (statuts, rôles, notifications) ;
    \item Chiffrement AES-256-GCM intégré nativement ;
    \item Portail public isolé pour les sociétés externes ;
    \item Coût nul en licences (stack open source) ;
    \item Déployable sur l'infrastructure existante de N+ONE.
\end{itemize}

\paragraph{Estimation des coûts d'infrastructure} En cas de déploiement en production,
le coût se limiterait à l'infrastructure serveur : environ 50 à 100 \euro/mois
sur un serveur cloud (ou nul si déployé sur les serveurs internes N+ONE), contre
plusieurs milliers d'euros annuels pour une licence SAP Ariba.


}
\end{spacing}

\newpage
